\documentclass[10pt,a4paper]{amsart}
\usepackage[margin=19mm]{geometry}
\usepackage{amsmath,amssymb,mathtools}
\usepackage[scr=rsfs,cal=euler]{mathalfa}
\usepackage{xfrac}
\usepackage[unicode,hidelinks,bookmarks=false]{hyperref}
\newcommand{\ie}{i.e.\ }
\newcommand{\cf}{cf.\ }
\newcommand{\resp}{respectively\ }
\newcommand{\Subsection}{\subsection}
\providecommand{\nobreakdash}{\nobreak\hbox{-}}
\setlength{\emergencystretch}{2em}
\multlinegap0pt
\usepackage[T1]{fontenc}
\usepackage{lmodern}
\catcode"2013=\active
\def–{\textendash}
\newcommand{\R}{\mathbb{R}}
\newcommand{\IR}{\mathbb{R}}
\newcommand{\Q}{\mathbb{Q}}
\newcommand{\IQ}{\mathbb{Q}}
\newcommand{\N}{\mathbb{N}}
\newcommand{\IN}{\mathbb{N}}
\newcommand{\Z}{\mathbb{Z}}
\newcommand{\IZ}{\mathbb{Z}}
\newcommand{\C}{\mathbb{C}}
\newcommand{\ind}{\mathds{1}}
\newcommand{\p}{\varphi}
\newcommand{\e}{\varepsilon}
\newcommand{\FF}{\mathscr{F}}
\newcommand{\dist}{\operatorname{dist}}
\newcommand{\rng}{\operatorname{rng}}
\newcommand{\supp}{\operatorname{supp}}
\renewcommand{\leq}{\leqslant}
\renewcommand{\geq}{\geqslant}
\newcommand{\mc}{\mathcal}
\newcommand{\on}{\operatorname}
\newcounter{maintheorem}
\renewcommand*{\themaintheorem}{\Alph{maintheorem}}
\newtheorem{mainth}[maintheorem]{Theorem}
\newcounter{maincorollary}
\renewcommand*{\themaincorollary}{\Alph{maincorollary}}
\newtheorem{maincor}[maincorollary]{Corollary}
\newtheorem{theorem}{Theorem}[section]
\newtheorem{lemma}[theorem]{Lemma}
\newtheorem{proposition}[theorem]{Proposition}
\newtheorem{corollary}[theorem]{Corollary}
\theoremstyle{definition}
\newtheorem{definition}[theorem]{Definition}
\newtheorem{question}[theorem]{Question}
\theoremstyle{remark}
\newtheorem{remark}{Remark}
\title[Filter $\pi$-bases]{The $\pi$-property of a Banach space along a filter: reduction to an analytic filter (Theorem B)}
\author{Tomasz Kania and Jaros{\l}aw Swaczyna}
\date{Source: arXiv:2508.21242v1 (CC BY 4.0)}
\begin{document}
\begin{abstract}
We examine the analyticity of the class of separable Banach spaces possessing the $\pi$-property, defined in terms of convergence along a filter. Our results establish that this class is $\Sigma^1_3$ whenever the underlying filter is analytic (as a subset of the Cantor set $\Delta$). Furthermore, we demonstrate that if the filter is countably generated, the class of such spaces is $\Sigma^1_2$ with respect to any admissible Polish topology on the family of closed subspaces of $C(\Delta)$.
\end{abstract}
\maketitle
\noindent\emph{Source and licence.} The $\pi$-property of a Banach space along a filter. Version arXiv:2508.21242v1 (CC BY 4.0), distributed under the Creative Commons Attribution 4.0 International licence, http://creativecommons.org/licenses/by/4.0/, as recorded in the arXiv OAI licence metadata for this exact version and shown on the abstract page https://arxiv.org/abs/2508.21242v1. Full rights notice: licenses/arxiv-2508.21242-CC-BY-4.0.txt.\par\smallskip
\noindent\textit{Editorial scope.} Counted unit: the second of the note's three main theorems, source label Th:B (source0.tex lines 176--178), printed as Theorem B in the published version and as Theorem B here, delivered together with its complete original author proof at source lines 660--684 as the single primary proof body. The proof is a detached block: the note states its three main theorems together in the Introduction and proves Theorem B much later, in the section headed Proof of Theorem A and Theorem B and Theorem C (source line 658); the block at source lines 660--684 opens with the author's own optional argument Proof of Theorem backslash-ref Th:B and is the only proof in the note that proves Theorem B, which is why the counted statement and the counted proof belong together. No mathematics is written, repaired or re-derived: the statement and the proof are the note's own text line for line. The note is already an amsart article (source0.tex line 2, class options 12pt and fleqn) that loads inputenc and fontenc with T1, so no publisher class is replaced or shipped: the wrapper is the lane's pinned amsart editorial wrapper, the note's own font encoding T1 and its own lmodern text font are reproduced in the wrapper preamble, and the note's own theorem environments (source lines 74--90, including the Alph-numbered mainth counter that prints Theorem A, Theorem B and Theorem C, the separate maincorollary counter, the section-numbered theorem counter shared by lemma, proposition, corollary, definition and question, and the definition and remark theorem styles) and its own macro block (source lines 45--71) are copied verbatim, so every printed number of the delivered unit is produced by the note's own counters and coincides with the published version: Section 1 Introduction, Definition 1.1, Question 1.2, Theorem A and the counted Theorem B. Required context retained uncounted: the whole Introduction of the note from its section heading to the line immediately preceding the counted statement (source lines 116--175), that is the motivation for replacing convergence in the strong operator topology by convergence along a filter, the definition of convergence along a filter, Definition 1.1 of the F-pi-property and of an F-pi-basis with the observation that the Frechet filter renders the familiar pi-property, the discussion of filter bases with continuous basis projections, of Kochanek's theorem for countably generated filters, of the results of the authors and of de Rancourt, Kania and Swaczyna, of Aviles-Kadets-Perez-Solecki on Baire filters and analytic filters that are not countably generated, the definition of a finite-dimensional decomposition, Question 1.2 asking whether the F-pi-property implies an FDD, the announcement of the complexity results with Theorem A, the paragraph explaining why Theorem B does not make higher projective complexities superfluous, and the lead-in sentence that announces the counted theorem. Every definition, theorem, question, remark, numbered display, footnote and citation of those source lines is retained. The note's abstract (source line 113) is reproduced verbatim in the wrapper front matter and the note's own title and author names are reproduced in the wrapper title block. Not part of this unit and not redistributed: the note's address, email, thanks, keywords and subject-class front matter (source lines 93--106); the statement of Theorem C (source lines 181--183); the whole of Sections 2 to 5 (source lines 184--744), that is the Preliminaries with the admissible topologies, the continuous-selection theorem, its Borel-selection corollary and the perturbation machinery for the proof of Theorem A, the proof of Theorem A itself, the proof of Theorem C, and the closing list of problems; and the note's bibliography listing, of which only the cited entries are transcribed. The delivered text cites seven keys (avilesetal, Casazza, RKS, GhaFund, HajekMedina, KaniaSwaczyna and Kochanek, eleven citation commands in all), and those seven entries are transcribed verbatim from the note's own in-source thebibliography (source lines 747--768) with the note's printed bibliography numbers preserved (avilesetal prints 2, Casazza 3, RKS 7, GhaFund 8, HajekMedina 12, KaniaSwaczyna 14 and Kochanek 16), so every printed citation number of the delivered unit is the published one. Reference adaptation: none was needed and none was made. The only labels and references inside the delivered unit are the note's own Definition 1.1 (label maindef, referred to in the Frechet-filter paragraph) and Question 1.2 (label q:1, referred to in the paragraph on the Borel class of the pi-property), together with the counted label Th:B itself, which the counted proof refers to in its own optional argument; all of them resolve inside the unit and no unresolved reference is produced. Printed slips kept literal: no printed word, symbol, hypothesis or number of the counted statement or of its proof was altered. The note's own forms are delivered as printed, among them the spelling of Frechet without its accent in the bibliography entries, the tag Th:B that names the statement rather than its printed letter, its own notation for filters on the natural numbers, its projection notation with subscripts, and its habit of writing the same inclusion both as a subset and as a superset relation in the construction of the generated filter. Non-ASCII characters: the source contains eleven characters outside ASCII in four distinct code points. Five of them are delivered: the accented e of the word Frechet at source lines 130, 133 (twice) and 172 inside the retained context, and the en dash of the page range 303 to 309 in the transcribed entry GhaFund at source line 755. The note itself loads inputenc with utf8 and fontenc with T1, and the wrapper reproduces both the T1 encoding and the note's own lmodern text font, so the accented e is delivered by the T1 mapping directly. The Unicode-to-T1 mapping used by XeTeX has no entry for the en dash, so the wrapper supplies the note's own inputenc mapping for that single code point: the character is made active and expands to the T1 en dash command, which is exactly what inputenc's utf8 option does in the note's own pdflatex run, so the delivered en dash is printed with the note's own glyph in the note's own font. No character is replaced, dropped or re-encoded, no mathematics is touched, and the xelatex receipt reports no missing character; no fidelity receipt for a character is needed. The note's loads of a4wide, scrextend, amssymb, amsmath, amsthm, mathtools, mathrsfs, euscript, dsfont, graphicx, enumerate, enumitem, accents, xcolor, todonotes and bbm are either already supplied by the wrapper's pinned preamble or unused by the retained spans, which contain no figure, no table, no backslash-includegraphics, no backslash-input and no to-do note; the wrapper's pinned mathalfa options set the calligraphic alphabet in Euler script and the script alphabet in rsfs where the note's own preamble used the default calligraphic alphabet and the euscript script alphabet, a purely typographic substitution of the wrapper that changes no mathematics, and the mathematics of the retained spans uses only the letters f, a, epsilon, alpha, sigma, theta, lambda, X and D in those alphabets. Layout: the wrapper's pinned class is amsart at 10pt on a4paper, while the note's own class options are 12pt and fleqn (source0.tex line 2), so the delivered unit is set at the wrapper's size and its displayed formulas are centred instead of flush left, and the note's own a4wide and scrextend page-layout loads are replaced by the wrapper's pinned geometry; these are page-layout differences of the wrapper only, no formula, symbol or argument changes. No mathematics is written, repaired or re-derived; all mathematics of the unit is native text with no image dependency.
\section{Introduction}
A Banach space is said to have the \(\pi\)-property if there exists a uniformly bounded net of finite-rank projections acting on the space that converges to the identity map in the strong operator topology, \emph{i.e.}, pointwise. For separable Banach spaces, the \(\pi\)-property is strongly connected to the existence of a finite-dimensional decomposition (FDD): the presence of an FDD ensures the \(\pi\)-property, yet whether the converse holds remains an open question (\cite[p.~3]{HajekMedina}). Notably, the stronger metric \(\pi\)-property does imply the existence of an FDD.

To further explore this relationship, we propose extending the framework to encompass a broader class of spaces. Specifically, we replace convergence in the strong operator topology with convergence along a filter, while maintaining the setting of the strong operator topology. This generalisation not only sheds new light on the \(\pi\)-property but also introduces a concept of intrinsic interest. To lay the groundwork for this study, we formally present the central definition underpinning this note that is expressed in terms of filter convergence. (Let \(\mathcal{F}\) be a filter on \(\N\), and let \(X\) be a Banach space. A sequence \((x_n)_{n=1}^\infty\) in \(X\) is said to converge to \(x \in X\) along \(\mathcal{F}\) if 
\[
    \lim_{n, \mathcal{F}} \|x_n - x\| = 0,
\]
meaning that for every \(\varepsilon > 0\), there exists \(A \in \mathcal{F}\) such that \(\|x_n - x\| < \varepsilon\) for all \(n \in A\).)

\begin{definition}[\(\mathcal{F}\)-\(\pi\)-property]\label{maindef}
Let \(\mathcal{F}\) be a filter on \(\mathbb{N}\). A (separable) Banach space \(X\) has the \(\mathcal{F}\)-\(\pi\)-\emph{property} if there exists a sequence of finite-rank projections \((P_n)_{n=1}^\infty\) on \(X\) such that, for every \(x \in X\),
\[
    \lim_{n, \mathcal{F}} \|P_n x - x\| = 0.
\]
In this case, the sequence \((P_n)_{n=1}^\infty\) is called an \(\mathcal{F}\)-\(\pi\)-\emph{basis} of \(X\). An \(\mathcal{F}\)-\(\pi\)-basis with respect to the Fréchet filter is simply referred to as a \(\pi\)-basis.
\end{definition}

For the sake of avoiding pathologies, we have stated the definition for filters on $\mathbb N$ thereby limiting ourselves to separable Banach spaces. Hereinafter, all filters considered contain the Fréchet filter, that is the filter of cofinite subsets of $\mathbb N$. It is readily seen that for the Fréchet filter, Definition~\ref{maindef} renders the familiar $\pi$-property. 

Banach spaces that have filter bases having continuous basis projections constitute paradigmatic examples of spaces with the $\FF$-$\pi$-property. Indeed, suppose that $(e_n)_{n=1}^\infty$ is an $\FF$-basis for a Banach space $X$. This means that for every $x\in X$ there exists a unique sequence of scalars $(e_n^*(x))_{n=1}^\infty$ such that
\begin{itemize}
    \item for every $n\in \mathbb N$, the assignment $y\mapsto e_n^*(y)$ is continuous,
    \item $\lim\limits_{n, \FF}\sum_{j=1}^n e_n^*(x)e_n =x$.
\end{itemize}
The finite-rank operators $(P_n)_{n=1}^\infty$ given by $P_n x = \sum_{j=1}^n e_n^*(x) e_n$ (basis projections) witness the $\FF$-$\pi$-property of $X$. \smallskip

Kochanek \cite{Kochanek} demonstrated that if \(\mathcal{F}\) is countably generated (or more generally, generated by fewer than \(\mathfrak{p}\) sets, where \(\mathfrak{p}\) denotes the pseudo-intersection number), that is, there exists a countable family (or a family of cardinality less than \(\mathfrak{p}\)) \(\mathscr{B} \subset \mathcal{F}\) such that for every \( A \in \mathcal{F} \), there exists \( B \in \mathscr{B} \) with \( B \subseteq A \), then the associated projections \((P_n)_{n=1}^\infty\) are uniformly bounded on some set in the filter, and the continuity assumption in the definition of an \(\mathcal{F}\)-basis becomes redundant. 

Moreover, the present authors have recently shown that the continuity assumption in the definition of an \(\mathcal{F}\)-basis is redundant for projective filters under proper set-theoretic assumptions (\cite{KaniaSwaczyna}) and even for analytic filters in ZFC (jointly with de Rancourt \cite{RKS}). 

The proof techniques from \cite{Kochanek} inspired Avil\'es \emph{et al.}~\cite{avilesetal} to identify the class of filters possibly generated by \(\mathfrak{c}\)-sized sets for which Kochanek's argument remains valid. A filter \(\mathcal{F}\) is termed \emph{Baire} if, for every complete metric space \( X \) and every family \(\{X_A \colon A \in \mathcal{F}\}\) of nowhere-dense subsets of \( X \) satisfying \( X_A \subseteq X_B \) whenever \( A, B \in \mathcal{F} \) and \( A \subseteq B \), it holds that \(\bigcup_{A \in \mathcal{F}} X_A \neq X\). 

Under Martin's Axiom, the classes of filters that are Baire with respect to all complete metric spaces and filters generated by fewer than continuum-many sets coincide. However, there exist models of ZFC that admit filters generated by exactly \(\mathfrak{c}\)-sized sets that are Baire (\cite[Section 5]{avilesetal}). On the other hand, analytic filters that are not countably generated fail to be Baire with respect to some Polish space (\cite[Proposition 4.1]{avilesetal}), which constrains the range of spaces with the \(\mathcal{F}\)-\(\mathcal{\pi}\)-property derived from \(\mathcal{F}\)-bases. \smallskip

An \emph{FDD} (finite-dimensional decomposition) for a Banach space \(X\) is a sequence of finite-dimensional subspaces \((F_n)_{n=1}^\infty\) such that every \(x \in X\) admits a unique representation
\[
    x = \sum_{n=1}^\infty x_n, \quad \text{with} \quad x_n \in F_n \ \text{for all} \ n \in \mathbb{N}.
\]
Let us formulate the following question, which encapsulates our motivation for investigating spaces with the \(\mathcal{F}\)-\(\mathcal{\pi}\)-property.

\begin{question}\label{q:1}
    Suppose that a separable Banach space has the \(\mathcal{F}\)-\(\mathcal{\pi}\)-property with respect to a certain filter \(\FF\). Does \( X \) have a finite-dimensional decomposition (FDD)?
\end{question}

The purpose of this note is to evaluate the complexity of spaces with the \(\mathcal{F}\)-\(\mathcal{\pi}\)-property within the framework of analytic filters. This investigation is inspired by a similar result of Ghawadrah \cite{GhaFund}, which concerns spaces with the usual \(\pi\)-property.

To state our result, we introduce the set \(\mathsf{SB}_{\mathcal{F}\text{-}\pi}\), a subset of \(\mathsf{SB}\) (formally defined in the following section) consisting of Banach spaces that exhibit the \(\mathcal{F}\)-\(\mathcal{\pi}\)-property.


\begin{mainth}\label{Th:A}Let $\FF$ be an analytic (or $\Sigma^1_{s}$) filter on $\mathbb N$. Then:
\begin{itemize}
    \item the set $\mathsf{SB}_{\mathcal{F}\text{-}\pi}$ is $\Sigma_3^1$ ($\Sigma_{s+2}^1$) in $\mathsf{SB}$.
    \item if $\FF$ is countably generated, the set $\mathsf{SB}_{\mathcal{F}\text{-}\pi}$ is $\Sigma^1_2$ in $\mathsf{SB}$.
\end{itemize}
\end{mainth}

For the Fréchet filter $\mathcal{F}_{{\rm Fr}}$, that is, in the case of the usual \(\pi\)-property, Ghawadrah~\cite{GhaFund} proved that $\mathsf{SB}_{\mathcal{F}_{{\rm Fr}}\text{-}\pi}$ is $\Sigma^0_6$. Even though neither Ghawadrah's identification of the Borel class of $\mathsf{SB}_{\mathcal{F}_{{\rm Fr}}\text{-}\pi}$ nor our identification of $\mathsf{SB}_{\mathcal{F}\text{-}\pi}$ is proved to be optimal, the increase in Borel complexity by one appears to be inevitable. This suggests a negative answer to Question~\ref{q:1}, should the $\pi$-property be equivalent to the existence of an FDD for separable Banach spaces (see~\cite[Theorem 6.4]{Casazza}).

Note also that the family $\mathsf{SB}_{\mathcal{F}\text{-}\pi}$ may depend on the particular choice of the filter $\mathcal{F}$. The following theorem asserts that a certain reduction to an analytic filter is possible.
%

\begin{mainth}\label{Th:B}
Let $\mathcal{F}$ be a filter on $\mathbb{N}$. Suppose that $X$ is a Banach space with the $\mathcal{F}$-$\pi$-property, and let $(P_n)_{n=1}^\infty$ witness this property. Then there exists an analytic filter $\mathcal{F}' \subset \mathcal{F}$ such that $(P_n)_{n=1}^\infty$ also witnesses the $\mathcal{F}'$-$\pi$-property of $X$.
\end{mainth}

\begin{proof}[Proof of Theorem~\ref{Th:B}]
Set 
\[
    \mc{A}:= \Bigg\{A \subset \IN\colon \exists_{x \in X} \exists_{\varepsilon>0}\; A \supset \Big\{ n \in \IN\colon \big\|P_nx - x \big\| \leq \varepsilon \Big\}\Bigg\}.
\]
Clearly $\mc{A} \subset \FF$. Note that $(P_n)$ witnesses that $X$ has $\FF^\prime$-$\pi$-property for any filter $\FF'$ on $\IN$ such that $\mc{A} \subset \FF' $. Now, for every $n \in \N$, consider the set
\[
\mc{B}_n= \Big\{ (x,\varepsilon,A)\in X \times (0,\infty)\times \{0, 1\}^\IN\colon n \in A \vee \big\|P_nx - x \big\| > \varepsilon \Big\}.
\]
By the continuity of $P_n$ we obtain that $\mc{B}_n$ is open. Observing that
\[
\mc{A} = \on{proj}_{\{0, 1\}^\IN}\Big[\bigcap_{n \in \IN} \mc{B}_n\Big],
\]
we deduce that $\mc{A}$ is analytic.\smallskip

As $\mc{A} \subset \FF$, finite intersections on elements of $\mc{A}$ are readily non-empty. Consequently, $\mc{A}$ generates a filter $\FF' \subset \FF$. We have:
\[
    \begin{aligned}
    \FF^\prime= &\{A \supset \IN\colon \exists_{n \in \IN} \exists_{A_1, A_2, \ldots, A_n \in \mc{A}} A \supset A_1 \cap A_2 \cap \ldots \cap A_n \} \\
    = & \bigcup_{n \in \IN} \on{proj}_{(\{0, 1\}^\N)_1} [\{ (A, A_1, \ldots, A_n)\in (\{0, 1\}^\N)^{n+1}\colon \\
    & A \supset A_1 \cap \ldots \cap A_n \wedge A_1, \ldots, A_n \in \mc{A}\}],
\end{aligned}
\]
where $\on{proj}_{(\{0, 1\}^\N)_1}$ is the projection onto the first coordinate. Consequently, $\FF'$ is an analytic filter and $(P_n)_{n=1}^\infty$ witnesses that $X$ has $\FF'$-$\pi$-property since $\mc{A}\subset \FF' \subset \FF$.
\end{proof}

\begin{thebibliography}{99}
\bibitem[2]{avilesetal} A. Avil\'es, V. Kadets, A. P\'erez, S. Solecki, Baire theorem for ideals of sets, \emph{J. Math. Anal. Appl.} \textbf{445} (2) (2017) 1221--1231.
\bibitem[3]{Casazza} P.~G. Casazza. \emph{Approximation properties}. In: Handbook of the geometry of Banach spaces, Vol. I. North-Holland, Amsterdam, 2001, pp. 271--316.
\bibitem[7]{RKS} N.~de~Rancourt, T.~Kania and J.~Swaczyna, Continuity of coordinate functionals of filter bases in Banach spaces. \emph{J. Funct. Anal.}. \textbf{284} (2023), no. 9, 109869.
\bibitem[8]{GhaFund} G. Ghawadrah, The descriptive complexity of approximation properties in an admissible topology. \emph{Fund. Math.} \textbf{249} (2020), no. 3, 303–309.
\bibitem[12]{HajekMedina} P.~H\'ajek and R.~Medina,
Compact retractions and Schauder decompositions in Banach spaces,
\emph{Trans. Amer. Math. Soc.} \textbf{376} (2023), 1343--1372.
\bibitem[14]{KaniaSwaczyna} T.~Kania, J.~Swaczyna, Large cardinals and continuity of coordinate functionals of filter bases in Banach spaces. \emph{Bull. Lond. Math. Soc}. \textbf{53} (2021), no. 1, 231--239.
\bibitem[16]{Kochanek} T.~Kochanek, $\mathcal{F}$-bases with brackets and with individual brackets in Banach spaces, \emph{Studia Math.} \textbf{211} (2012), 259--268.
\end{thebibliography}
\end{document}
